\section{Parte A}

\definecolor{headercolor}{HTML}{C2185B}
\definecolor{rowcolor}{HTML}{FCE4EC}

\title{Prueba}
\author{Candelaria Grigera}
\date{April 2026}

\section{Primer Experiencia}


\subsection{Instrumentos y componentes utilizados}

\begin{itemize}
    \item Fuente de tensión alterna $12\ V - 50\ Hz$.
    \item Resistencia de $68\ \Omega$.
    \item Bobina de 150 vueltas con resistencia desconocida.
    \item Multímetros digitales
    \item Cables banana-banana para el conexionado.
\end{itemize}

\subsection{Circuito implementado}

El circuito implementado para la primera experiencia fue el presentado en la fig. \ref{fig:circuito1}

\begin{figure}[H]
  \centering
  \begin{circuitikz}[american, scale=1.0, transform shape]
  
    \coordinate (A) at (0, 0);
    \coordinate (B) at (3, 0);
    \coordinate (G) at (8, 0);
    \coordinate (D) at (10, 0);
    \coordinate (E) at (10, -3);
    \coordinate (F) at (0, -3);
   
    \draw (A) to[R, l=$R$] (B);
    \draw (B) to[L, l=$L$] (5, 0);
    \draw (5,0) to[R, l=$r$] (G);
    \draw (G) to[rmeter, t=A] (D);
   
    \draw (E) -- (12,-3);
    \draw (12,-3) to[sV] (12,0);
    \node[right=4pt] at (12.5,-1.3) {$12\,\text{V}$};
    \node[right=4pt] at (12.5,-1.8) {$50\,\text{Hz}$};
    \draw (12,0) -- (D);
   
    \draw (F) -- (E);
    \draw (A) -- (F);
   
    \fill (B) circle (2.5pt);
    \fill (G) circle (2.5pt);
   
    \coordinate (VtL) at ($(A)+(0,1.6)$);
    \coordinate (VtR) at ($(G)+(0,1.6)$);
    \draw[dotted] (A) -- (VtL);
    \draw[dotted] (G) -- (VtR);
    \draw[dotted] (VtL) -- (VtR);
    \node[draw, circle, fill=white, inner sep=4pt] at ($(VtL)!0.5!(VtR)$) {V};
   
    \coordinate (VRL) at ($(A)+(0,-1.6)$);
    \coordinate (VRR) at ($(B)+(0,-1.6)$);
    \draw[dotted] (A) -- (VRL);
    \draw[dotted] (B) -- (VRR);
    \draw[dotted] (VRL) -- (VRR);
    \node[draw, circle, fill=white, inner sep=4pt] at ($(VRL)!0.5!(VRR)$) {$V_R$};
   
    \coordinate (VZLL) at ($(B)+(0,-1.6)$);
    \coordinate (VZLR) at ($(G)+(0,-1.6)$);
    \draw[dotted] (B) -- (VZLL);
    \draw[dotted] (G) -- (VZLR);
    \draw[dotted] (VZLL) -- (VZLR);
    \node[draw, circle, fill=white, inner sep=4pt] at ($(VZLL)!0.5!(VZLR)$) {$V_{ZL}$};
   
    \draw[dashed] (2.5,-0.7) rectangle (8.5,0.7);
    \node[above] at (5.5,0.7) {$Z_L$};

  \end{circuitikz}
  \caption{Circuito 1}
  \label{fig:circuito1}
\end{figure}

\subsection{Objetivos}

Determinar de manera experimental la resistencia interna y la inductancia de la bobina proporcionada por la cátedra.

\subsubsection{Datos Obtenidos}

Los datos que se obtuvieron a partir de la medición fueron los presentados en la tabla \ref{tab:mediciones1}

\begin{table}[H]
  \centering
  \renewcommand{\arraystretch}{1.4}
  \setlength{\arrayrulewidth}{0.5pt}
  \resizebox{\columnwidth}{!}{%
  \begin{tabular}{|
    >{\centering\arraybackslash}m{1.5cm}|
    >{\centering\arraybackslash}m{1.5cm}|
    >{\centering\arraybackslash}m{1.5cm}|
    >{\centering\arraybackslash}m{1.8cm}|
    >{\centering\arraybackslash}m{1.8cm}|
    >{\centering\arraybackslash}m{1.8cm}|
  }
  \hline
  \rowcolor{headercolor}
  \textcolor{white}{\textbf{$I$ [mA]}} &
  \textcolor{white}{\textbf{$V_f$ [V]}} &
  \textcolor{white}{\textbf{$V_R$ [V]}} &
  \textcolor{white}{\textbf{$V_{ZL}$ [V]}} &
  \textcolor{white}{\textbf{$R$ [$\Omega$]}} &
  \textcolor{white}{\textbf{$r$ [$\Omega$]}} &
  \hline
  \rowcolor{rowcolor}
  107 & 12,1 & 5,1 & 8,4 & 52 & 48 \\
  \hline
  \end{tabular}}
  \caption{Mediciones experiencia 1}
  \label{tab:mediciones1}
\end{table}
Cabe destacar que dichos valores fueron medidos a partir de los multímetros utilizados. 

A partir de los datos obtenidos y utilizando la siguiente ecuación:

\begin{equation*}
    \varphi = \cos^{-1}\!\left(\frac{V_R^2 + V_F^2 - V_{ZL}^2}{2V_R V_F}\right)
\end{equation*}

Se obtiene que $\varphi=34.38^{\circ}$.

También podemos obtener a partir de los datos y del valor de $\varphi$ obtenido tanto la tensión sobre la resistencia interna de la bobina como la tensión sobre la inductancia de la bobina. Dichos valores se obtuvieron a partir de las siguientes fórmulas:

\begin{align*}
  V_L &= V_F \sin(\varphi) \\
  V_r &= V_F \cos(\varphi) - V_R
\end{align*}

Obteniéndose así:

\begin{align*}
    V_L &=  \\
    V_r &=
\end{align*}

\textbf{NO TERMINO DE ENTENDER LA TABLA DEL SHEET Y NO ME CIERRAN LAS CUENTAS}





\newpage